\section{Decoding a rational reflection on finite words}
\label{sec:decoder}

The matrix identities for a standard directive supply many pairs related
by its reflection: one may exchange the blocks $ab$ and $ba$ in any
concatenation. The question is whether there are other pairs. We prove
that there are none, even among arbitrary finite words. The proof reads
two letters at a time. A change of projective coordinate places all
possible first letters in two fixed intervals, and the final letter of
the directive determines which inequalities exclude a wrong continuation.
\Cref{subsec:centre-decoder} proves the same statement for the reflection of
every palindrome by a shorter argument in the slope coordinate. The cylinders
of this section are kept because they also measure how far a reflected column
can move, which is what \cref{sec:displaced} needs.

For a directive $d$, retain the notation
$$
 U=\mu(\phi_d(a)),\qquad V=\mu(\phi_d(b)),\qquad
 G=\mu(\operatorname{Pal}(d)),\qquad
 \nu_d(w)=\mu(\phi_d(w)).
$$
Thus $\nu_d$ evaluates a word by replacing $a,b$ with $U,V$.
Write $M_0=e^{\mathsf T}Ge$, $\rho_0=(Ge)_1$ and
$$
 K=2ee^{\mathsf T}G-M_0I,\qquad R=K/M_0.
$$
Since $(ee^{\mathsf T}G)^2=M_0ee^{\mathsf T}G$, one has
$K^2=M_0^2I$; in particular $R$ is an invertible involution.
For a word $x$ in the two blocks $ab,ba$, let $\widehat{x}$ exchange
these blocks in their original order. The parsing into blocks is
unambiguous because both blocks have length two. For example,
$\widehat{abba}=baab$. For a nonzero column $v$, the notation $[v]$
denotes its line in the real projective line.

\begin{theorem}[Finite-word reflection decoder]\label{thm:word-decoder}
For every finite directive $d$ and every pair of finite words $x,y$,
including the empty word,
\begin{equation}\label{eq:decoder-projective}
 [R\nu_d(x)e]=[\nu_d(y)e]
 \quad\Longleftrightarrow\quad
 x\in\{ab,ba\}^*,\qquad y=\widehat{x}.
\end{equation}
Whenever these conditions hold, the proportionality is exact:
\begin{equation}\label{eq:decoder-exact}
 R\nu_d(x)e=\nu_d(\widehat{x})e.
\end{equation}
\end{theorem}

There is no palindrome or primitive-count hypothesis in this theorem.
Those restrictions will enter only when the decoder is applied to the
midpoint problem.

\subsection{The two block identities}

The literal standard identities proved in the preceding section give
\begin{equation}\label{eq:decoder-standard}
 UV=ABG,\qquad VU=BAG,\qquad G=G^{\mathsf T}.
\end{equation}
Appending $a$ to $d$ replaces $(U,V)$ by $(U,VU)$, and appending
$b$ replaces it by $(UV,V)$. The direction of this update will matter
when we compare the two generators.

Direct multiplication gives $AB+BA=6ee^{\mathsf T}$. For every
symmetric matrix $G$ one also has
$\tr(ABG)=3e^{\mathsf T}Ge$. Consequently
\begin{equation}\label{eq:decoder-intertwiners}
 3K=UV+VU-\tr(UV)I,\qquad
 KUV=VUK,\qquad KVU=UVK,\qquad Ke=M_0e.
\end{equation}
Here the middle identities can be checked without a word argument:
Cayley--Hamilton makes
$$
 (UV+VU-\tr(UV)I)UV
 =VUUV-I
 =VU(UV+VU-\tr(UV)I).
$$
Interchanging $U,V$ proves the other identity. In particular, these
relations prove the forward construction in
\eqref{eq:decoder-exact} by induction on the number of blocks. The
rest of the section proves exhaustion.

\subsection{A common coordinate at the terminal slope}

Matrices act on slopes by
$$
 \begin{pmatrix}a&b\\c&d\end{pmatrix}(t)
 =\frac{at+b}{ct+d},
$$
with the usual interpretation at infinity. The terminal column $e$
has slope two. Set
\begin{equation}\label{eq:decoder-interval}
 \mathcal J=[8/5,29/12],\qquad z=\frac1{t-2}.
\end{equation}
Conjugation by $P=\left(\begin{smallmatrix}2&1\\1&0\end{smallmatrix}\right)$
is exactly this change of coordinate, and sends the terminal slope to
$z=\infty$. The usefulness of the coordinate is that every positive
word already lies in the same interval:
$$
 A(\mathcal J)=[21/13,70/41]\subset\operatorname{int}\mathcal J,\qquad
 B(\mathcal J)=[50/21,169/70]\subset\operatorname{int}\mathcal J.
$$
Both matrices have determinant one, so both maps are increasing on
$\mathcal J$ and these endpoint evaluations prove the inclusions. The
second is the tight one: its upper endpoint clears that of $\mathcal J$ by
$29/12-169/70=1/420$, so the inclusion must be read in exact arithmetic.
Since $2\in\mathcal J$, induction applies to every finite positive word.
The interval $\mathcal J$ is a rational neighbourhood of the interval between the
attracting fixed points of $A$ and $B$ on slopes, the golden ratio $(1+\sqrt5)/2$ and
$1+\sqrt2$; its endpoints $8/5$ and $29/12$ are a Fibonacci and a Pell quotient just
outside them. Moreover $A(t)<2<B(t)$ for every positive $t$, so the first letter of a
nonempty word decides on which side of the terminal slope its column lies. In the
coordinate $z$ the two halves $[8/5,2)$ and $(2,29/12]$ of $\mathcal J$ become the rays
$z\leqslant-5/2$ and $z\geqslant12/5$, and the terminal slope $2$ becomes $z=\infty$.

There are positive integers $p,q$ such that
\begin{equation}\label{eq:decoder-trace-column}
 \tr U=3p,\quad (P^{-1}UP)_{21}=-p,\qquad
 \tr V=3q,\quad (P^{-1}VP)_{21}=q.
\end{equation}
At the empty directive this holds with $(p,q)=(1,2)$.
For the induction step, multiplication with a general symmetric $G$
gives
$$
 \tr(ABG)=\tr(BAG)=3M_0,\qquad
 (P^{-1}ABGP)_{21}=-M_0,\qquad
 (P^{-1}BAGP)_{21}=M_0.
$$
Thus the two directive updates replace $q$ or $p$, respectively,
by the positive integer $M_0$.

Define the positive reciprocal parameters and two coordinate centers by
$$
 \alpha_0=p^{-2},\qquad\beta_0=q^{-2},\qquad
 u=\frac{(P^{-1}UP)_{11}}p,\qquad
 v=\frac{(P^{-1}VP)_{11}}q.
$$
Trace and determinant determine the remaining entries:
\begin{equation}\label{eq:decoder-cusp-matrices}
 \begin{split}
 P^{-1}UP&=
 p\begin{pmatrix}u&u^2-3u+\alpha_0\\-1&3-u\end{pmatrix},\\
 P^{-1}VP&=
 q\begin{pmatrix}v&3v-v^2-\beta_0\\1&3-v\end{pmatrix}.
 \end{split}
\end{equation}
Their projective maps in the new coordinate are therefore
\begin{equation}\label{eq:decoder-fg}
 f(z)=-u-\frac{\alpha_0}{z+u-3},\qquad
 g(z)=v-\frac{\beta_0}{z+3-v}.
\end{equation}

\begin{lemma}[Uniform bounds and the last directive letter]
\label{lem:decoder-dominance}
For every directive,
\begin{equation}\label{eq:decoder-uv}
 3\leqslant u\leqslant41/12,\qquad
 5/2\leqslant v\leqslant21/8.
\end{equation}
If $L=u+v-3$, then
\begin{equation}\label{eq:decoder-curve}
 M_0=pqL,\qquad
 \alpha_0+\beta_0=L(3-L),\qquad 5/2\leqslant L<3.
\end{equation}
In particular $p\geqslant1$ and $q\geqslant2$. For a nonempty
directive the following alternatives hold:
\begin{equation}\label{eq:decoder-dominance}
 \begin{array}{c@{\quad}c@{\quad}l}
 \text{last letter}&\text{trace dominance}&\text{reciprocal bounds}\\[2pt]
 a&q\geqslant5p&\beta_0\leqslant\alpha_0/25,\quad\alpha_0\leqslant1\\
 b&p\geqslant(5/2)q&\alpha_0\leqslant4\beta_0/25,\quad\beta_0\leqslant1/4.
 \end{array}
\end{equation}
\end{lemma}

\begin{proof}
The word $\phi_d(a)$ is either $a$ or begins with $ab$; the word
$\phi_d(b)$ is either $b$ or begins with $ba$. This follows by
induction from the two updates following \eqref{eq:decoder-standard}.
Consequently $U(2)=A(s)$, where $s=2$ if $U=A$ and otherwise $s$ is the slope of a column
of a word beginning with $b$, so that $s\in B(\mathcal J)\subset[2,29/12]$. Likewise
$V(2)=B(s')$ with $s'=2$ or $s'\in A(\mathcal J)\subset[8/5,2]$. Since $A$ and $B$ are
increasing,
$$
 5/3=A(2)\leqslant U(2)\leqslant A(29/12)=70/41,\qquad
 50/21=B(8/5)\leqslant V(2)\leqslant B(2)=12/5.
$$
On the other hand \eqref{eq:decoder-cusp-matrices} gives $U(2)=2-1/u$ and
$V(2)=2+1/v$: the column $e$ has cusp coordinate $\infty$, and $f(\infty)=-u$,
$g(\infty)=v$ by \eqref{eq:decoder-fg}. Solving, $12/41\leqslant1/u\leqslant1/3$ and
$8/21\leqslant1/v\leqslant2/5$, which is \eqref{eq:decoder-uv}.

The lower-left entry of $P^{-1}VUP$ is $pq(u+v-3)$, so it
equals $M_0=pqL$. The matrix
$P^{-1}(UV+VU)P=6P^{-1}ee^{\mathsf T}GP$ has zero second row.
Its lower-right entry, expanded with
\eqref{eq:decoder-cusp-matrices}, gives
$\alpha_0+\beta_0=L(3-L)$. The bounds on $u,v$ give $L\geqslant5/2$;
positivity of $\alpha_0+\beta_0$ then gives $L<3$.

This argument has not required lower bounds on $p,q$ beyond their
positive integrality. It now shows that either update to $pqL$
strictly increases the replaced coordinate. Starting at $(1,2)$,
we obtain $p\geqslant1,q\geqslant2$ at every stage.
For an appended $a$, the new $q$ is
$L_{\rm old}p q_{\rm old}\geqslant5p$.
For an appended $b$, the new $p$ is
$L_{\rm old}p_{\rm old}q\geqslant(5/2)q$.
Taking reciprocal squares gives \eqref{eq:decoder-dominance}.
\end{proof}

\subsection{First-letter intervals and the reflection}

Let $\mathcal C$ be the set of cusp slopes of all $\nu_d(w)e$.
The empty word contributes infinity, and
\begin{equation}\label{eq:decoder-orbit}
 \mathcal C=\{\infty\}\cup f(\mathcal C)\cup g(\mathcal C)
 \subset\mathcal H
 :=\{z\leqslant-5/2\}\cup\{z\geqslant12/5\}\cup\{\infty\}.
\end{equation}
The inclusion follows by transforming $\mathcal J$ under
\eqref{eq:decoder-interval}, as explained above. Each of the maps \eqref{eq:decoder-fg} has
a single pole: $f$ at $z=3-u$ and $g$ at $z=v-3$. By \eqref{eq:decoder-uv} the first lies
in $[-5/12,0]$ and the second in $[-1/2,-3/8]$, so both poles sit in the gap between the
two rays of $\mathcal H$, and their distances to the rays are bounded below as in
\cref{tab:poles}. On the positive ray the denominators $z+u-3$ and $z+3-v$ are these
distances, and on the negative ray they are their negatives; in particular no denominator
vanishes on $\mathcal H$.

\begin{table}[htbp]
\centering
\caption{The poles of the two first-letter maps \eqref{eq:decoder-fg} and lower bounds for
their distances to the two rays of $\mathcal H$, from \eqref{eq:decoder-uv}. Dividing the
numerators $\alpha_0$ and $\beta_0$ by these distances gives the distances from the centres
$-u$ and $v$ to the four endpoints of the cylinders \eqref{eq:decoder-cylinders}.}
\label{tab:poles}
\small
\renewcommand{\arraystretch}{1.15}
\begin{tabular}{@{}lccc@{}}
\toprule
map & pole & distance to $z\geqslant12/5$ & distance to $z\leqslant-5/2$\\
\midrule
$f$ (letter $a$) & $3-u\in[-5/12,\,0]$ & $\geqslant12/5-0=12/5$ & $\geqslant-5/12+5/2=25/12$\\
$g$ (letter $b$) & $v-3\in[-1/2,\,-3/8]$ & $\geqslant12/5+3/8=111/40$ & $\geqslant-1/2+5/2=2$\\
\bottomrule
\end{tabular}
\end{table}

A map $z\mapsto c-\beta/(z-z_0)$ with $\beta>0$ sends every point at distance at least
$\delta$ from its pole $z_0$ to within $\beta/\delta$ of $c$, below $c$ to the right of the
pole and above $c$ to the left of it, and it sends $\infty$ to $c$. For $f$ the centre is
$c=-u$, the numerator $\alpha_0$ and the pole $3-u$; for $g$ they are $v$, $\beta_0$ and
$v-3$. Applied to the four entries of \cref{tab:poles}, this gives
\begin{equation}\label{eq:decoder-cylinders}
 \begin{split}
 f(\mathcal H)&\subset
 \mathcal U:=[-u-5\alpha_0/12,-u+12\alpha_0/25],\\
 g(\mathcal H)&\subset
 \mathcal V:=[v-40\beta_0/111,v+\beta_0/2].
 \end{split}
\end{equation}
For example, the lower end of $\mathcal V$ is $v-\beta_0\cdot40/111$: the point of the
positive ray closest to the pole of $g$ is $12/5$, at distance $12/5-(v-3)\geqslant111/40$.
\Cref{fig:cusp-cylinders} shows the configuration at the empty directive.

\begin{figure}[htbp]
\centering
% Author: Dr. Denys Dutykh (Khalifa University of Science and Technology, Abu Dhabi, UAE)
% Generated by supplement/code/make_figures.py from exact data; do not edit by hand.

\begin{tikzpicture}[font=\footnotesize, >={Stealth[length=4pt]}]
\draw[mist, densely dotted, line width=0.5pt] (2.871,-0.56) -- (2.871,-1.54);
\draw[mist, densely dotted, line width=0.5pt] (12.573,-0.56) -- (12.573,-1.54);
\draw[mist, line width=0.8pt] (2.871,0) -- (12.573,0);
\draw[navy, line width=2.2pt] (2.871,0) -- (0.15,0);
\draw[navy, line width=2.2pt] (12.573,0) -- (13.71,0);
\draw[navy, ->, line width=0.8pt] (0.2,0) -- (-0.05,0);
\draw[navy, ->, line width=0.8pt] (13.66,0) -- (13.91,0);
\node[ink, below] at (0,-0.08) {$\infty$};
\node[ink, below] at (13.86,-0.08) {$\infty$};
\node[ink, below] at (2.871,-0.08) {$-5/2$};
\node[ink, below] at (12.573,-0.08) {$12/5$};
\node[ink, below] at (7.821,-0.08) {$0$};
\draw[ink, line width=0.5pt] (7.821,-0.07) -- (7.821,0.07);
\filldraw[fill=sand, draw=ink, line width=0.4pt] (1.056,0.16) rectangle (2.831,0.58);
\node[ink, above] at (1.944,0.6) {$\mathcal U$};
\filldraw[fill=sand, draw=ink, line width=0.4pt] (12.593,0.16) rectangle (13.018,0.58);
\node[ink, above] at (12.806,0.6) {$\mathcal V$};
\draw[navy, line width=0.25pt] (1.061,0.2) -- (1.061,0.54);
\draw[navy, line width=0.25pt] (1.062,0.2) -- (1.062,0.54);
\draw[navy, line width=0.25pt] (1.063,0.2) -- (1.063,0.54);
\draw[navy, line width=0.25pt] (1.089,0.2) -- (1.089,0.54);
\draw[navy, line width=0.25pt] (1.115,0.2) -- (1.115,0.54);
\draw[navy, line width=0.25pt] (1.116,0.2) -- (1.116,0.54);
\draw[navy, line width=0.25pt] (1.119,0.2) -- (1.119,0.54);
\draw[navy, line width=0.25pt] (1.123,0.2) -- (1.123,0.54);
\draw[navy, line width=0.25pt] (1.124,0.2) -- (1.124,0.54);
\draw[navy, line width=0.25pt] (1.125,0.2) -- (1.125,0.54);
\draw[navy, line width=0.25pt] (1.881,0.2) -- (1.881,0.54);
\draw[navy, line width=0.25pt] (2.461,0.2) -- (2.461,0.54);
\draw[navy, line width=0.25pt] (2.463,0.2) -- (2.463,0.54);
\draw[navy, line width=0.25pt] (2.466,0.2) -- (2.466,0.54);
\draw[navy, line width=0.25pt] (2.541,0.2) -- (2.541,0.54);
\draw[navy, line width=0.25pt] (2.612,0.2) -- (2.612,0.54);
\draw[navy, line width=0.25pt] (2.613,0.2) -- (2.613,0.54);
\draw[navy, line width=0.25pt] (2.624,0.2) -- (2.624,0.54);
\draw[navy, line width=0.25pt] (2.634,0.2) -- (2.634,0.54);
\draw[navy, line width=0.25pt] (2.635,0.2) -- (2.635,0.54);
\draw[navy, line width=0.25pt] (2.637,0.2) -- (2.637,0.54);
\draw[navy, line width=0.25pt] (12.601,0.2) -- (12.601,0.54);
\draw[navy, line width=0.25pt] (12.606,0.2) -- (12.606,0.54);
\draw[navy, line width=0.25pt] (12.611,0.2) -- (12.611,0.54);
\draw[navy, line width=0.25pt] (12.612,0.2) -- (12.612,0.54);
\draw[navy, line width=0.25pt] (12.771,0.2) -- (12.771,0.54);
\draw[navy, line width=0.25pt] (12.941,0.2) -- (12.941,0.54);
\draw[navy, line width=0.25pt] (12.942,0.2) -- (12.942,0.54);
\draw[navy, line width=0.25pt] (12.943,0.2) -- (12.943,0.54);
\draw[navy, line width=0.25pt] (12.969,0.2) -- (12.969,0.54);
\draw[navy, line width=0.25pt] (12.995,0.2) -- (12.995,0.54);
\draw[navy, line width=0.25pt] (12.996,0.2) -- (12.996,0.54);
\draw[navy, line width=0.25pt] (12.999,0.2) -- (12.999,0.54);
\draw[navy, line width=0.25pt] (13.003,0.2) -- (13.003,0.54);
\draw[navy, line width=0.25pt] (13.004,0.2) -- (13.004,0.54);
\draw[navy, line width=0.25pt] (13.005,0.2) -- (13.005,0.54);
\draw[wine, <->, line width=0.6pt] (2.871,-0.78) -- (6.996,-0.78);
\draw[wine, <->, line width=0.6pt] (7.821,-0.78) -- (12.573,-0.78);
\draw[wine, line width=3pt] (6.996,-0.78) -- (7.821,-0.78);
\node[wine, above] at (4.934,-0.78) {$25/12$};
\node[wine, above] at (10.197,-0.78) {$12/5$};
\node[wine, below] at (7.409,-0.8200000000000001) {pole of $f$};
\draw[teal, <->, line width=0.6pt] (2.871,-1.42) -- (6.831,-1.42);
\draw[teal, <->, line width=0.6pt] (7.079,-1.42) -- (12.573,-1.42);
\draw[teal, line width=3pt] (6.831,-1.42) -- (7.079,-1.42);
\node[teal, above] at (4.851,-1.42) {$2$};
\node[teal, above] at (9.826,-1.42) {$111/40$};
\node[teal, below] at (6.955,-1.46) {pole of $g$};
\end{tikzpicture}

\caption{The cusp coordinate $z$ at the empty directive, where $f(z)=-3-1/z$ and
$g(z)=5/2-1/(4z+2)$. Every column of a nonempty word lies on one of the two rays of
$\mathcal H$ (thick), and the terminal column is the point at infinity. The poles of $f$ and
$g$ lie in the gap between the rays; the bars below the axis show the intervals $[-5/12,0]$ and
$[-1/2,-3/8]$ that contain them for every directive, and the arrows show the four distances of
\cref{tab:poles}, measured from the ends of these ranges. The shaded boxes are the cylinders
$\mathcal U$ and $\mathcal V$ of \eqref{eq:decoder-cylinders}, and the ticks above the axis are
the columns of all words of length at most nine.}
\label{fig:cusp-cylinders}
\end{figure}

These two finite intervals lie strictly in the negative and positive components of
$\mathcal H$:
$$
 -u+12\alpha_0/25\leqslant-3+12/25<-5/2,\qquad
 v-40\beta_0/111\geqslant5/2-10/111>12/5.
$$
Consequently a finite slope in $\mathcal C$ determines the first
letter, and inversion of its projective generator determines the
remainder. The terminal value cannot belong to either finite
interval. Induction therefore proves injectivity of finite-word
evaluation at the terminal slope.

Put $\tau=2\rho_0/M_0$. Since both entries of $Ge$ are positive,
$M_0=2\rho_0+(Ge)_2$ gives $0<\tau<1$. Direct conjugation and
\eqref{eq:decoder-intertwiners} give
\begin{equation}\label{eq:decoder-k}
 P^{-1}RP=\begin{pmatrix}1&\tau\\0&-1\end{pmatrix},\qquad
 k(z)=-z-\tau,\qquad
 \tau=u-v+\delta,\quad \delta=\frac{\alpha_0-\beta_0}{L}.
\end{equation}
For the last equality, expand the upper-right entry of
$3P^{-1}KP$ using \eqref{eq:decoder-cusp-matrices}.
The map $k$ fixes infinity and sends the two finite components
of $\mathcal H$ to opposite signs. It need not preserve
$\mathcal H$; whenever both endpoints belong to $\mathcal C$,
their first letters must nevertheless be opposite.

We will also use positivity after an intermediate reflection, even
when that reflected point has not yet been shown to be in the orbit.
In the original coordinate,
\begin{equation}\label{eq:decoder-positive-reflection}
 R(t)=2-\frac{t-2}{1+\tau(t-2)}>0\qquad(t\in\mathcal J).
\end{equation}
Indeed the denominator exceeds $3/5$. If $t<2$, the value exceeds
two; if $t\geqslant2$, it is at least $2-(t-2)\geqslant19/12$.
As $U$ begins with $A$ and $V$ with $B$, they send every positive
input below and above two, respectively. Hence $fk$ has negative
cusp slope on $\mathcal C$, while $gk$ has positive cusp slope.
These statements include the terminal input.

\subsection{Excluding a wrong second letter}

After opposite first letters have been removed, the two residual maps
are
\begin{equation}\label{eq:decoder-residuals}
 \begin{split}
 \mathcal E(z):=g^{-1}kf(z)
 &=\frac{\beta_0}{\delta-\alpha_0/(z+u-3)}+v-3,\\
 \mathcal D(z):=f^{-1}kg(z)
 &=\frac{\alpha_0}{\delta-\beta_0/(z+3-v)}+3-u.
 \end{split}
\end{equation}
These formulas follow by solving \eqref{eq:decoder-fg} for its
input. The identities \eqref{eq:decoder-intertwiners} imply
\begin{equation}\label{eq:decoder-residual-return}
 \mathcal E g=fk,\qquad \mathcal D f=gk.
\end{equation}
It remains to exclude an empty remainder or the wrong next letter.
We prove precisely
\begin{equation}\label{eq:decoder-wrong}
 \mathcal E(\{\infty\}\cup\mathcal U)\cap\mathcal C=\varnothing,
 \qquad
 \mathcal D(\{\infty\}\cup\mathcal V)\cap\mathcal C=\varnothing.
\end{equation}
Only these domains are needed. In particular, no estimate across a
pole of either residual map will be used.

\smallskip\noindent\emph{A directive ending in $a$.}
Here $\beta_0\leqslant\alpha_0/25$ and $\alpha_0\leqslant1$.
Since $L<3$,
$$
 \delta>\frac{8\alpha_0}{25}>0,\qquad
 0<\frac{\beta_0}{\delta}<\frac18.
$$
It follows from $-1/2\leqslant v-3\leqslant-3/8$ that
$\mathcal E(\infty)\in(-1/2,-1/4)$.
For $z\in\mathcal U$, the quantity $z+u-3$ is negative,
so the denominator of $\mathcal E$ is larger than $\delta$.
Therefore
\begin{equation}\label{eq:decoder-a-E}
 \mathcal E(\{\infty\}\cup\mathcal U)\subset(-1/2,-1/4),
\end{equation}
which is disjoint from $\mathcal H$.

For the other residual, direct simplification gives
$$
 \mathcal D(\infty)-v=\frac{L\beta_0}{\alpha_0-\beta_0}
 >\frac52\beta_0>\frac{\beta_0}{2}.
$$
If $z\in\mathcal V$, put $s=z+3-v$. The cylinder bound gives
$s\geqslant3-40\beta_0/111>2$. Its residual denominator $h_0$ satisfies
$$
 h_0:=\delta-\frac{\beta_0}{s}
 >\left(\frac8{25}-\frac1{50}\right)\alpha_0
 =\frac{3\alpha_0}{10}>0,\qquad h_0<\delta<\frac{\alpha_0}{L}.
$$
Using $L\delta=\alpha_0-\beta_0$ now yields
\begin{equation}\label{eq:decoder-a-D}
 \mathcal D(z)-v
 =\frac{\beta_0(1+L/s)}{h_0}
 >\frac{L\beta_0}{\alpha_0}\geqslant\frac52\beta_0>\frac{\beta_0}{2}.
\end{equation}
Both the terminal value and these finite values are positive and
strictly above the entire positive cylinder $\mathcal V$. They
cannot be in $\mathcal C$.

\smallskip\noindent\emph{A directive ending in $b$.}
Now $\alpha_0\leqslant4\beta_0/25$ and $\beta_0\leqslant1/4$. Set
$\eta=-\delta=(\beta_0-\alpha_0)/L$. Then
$$
 \eta>\frac{7\beta_0}{25}>0,\qquad
 0<\frac{\alpha_0}{\eta}<\frac47.
$$
At the terminal input,
$$
 \mathcal E(\infty)+u=-\frac{L\alpha_0}{\beta_0-\alpha_0}
 <-10\alpha_0.
$$
For $z\in\mathcal U$, put $s=3-z-u$.
Since $\alpha_0\leqslant1/25$, the cylinder bound gives
$s\geqslant3-12\alpha_0/25>2$.
The positive denominator
$$
 h_0:=\eta-\frac{\alpha_0}{s}
 >\left(\frac7{25}-\frac2{25}\right)\beta_0
 =\frac{\beta_0}{5},\qquad h_0<\eta<\frac{\beta_0}{L},
$$
allows the exact simplification
\begin{equation}\label{eq:decoder-b-E}
 \mathcal E(z)+u
 =-\frac{\alpha_0(1+L/s)}{h_0}
 <-\frac{L\alpha_0}{\beta_0}\leqslant-10\alpha_0.
\end{equation}
These values and the terminal value are negative and strictly below
$-u-5\alpha_0/12$, the lower endpoint of $\mathcal U$.
Thus they are outside $\mathcal C$.

Finally $\mathcal D(\infty)=-\alpha_0/\eta+3-u$.
For $z\in\mathcal V$, the quantity $s=z+3-v$ is positive and
$$
 \mathcal D(z)=-\frac{\alpha_0}{\eta+\beta_0/s}+3-u.
$$
The bounds $0\leqslant u-3\leqslant5/12$ and
$0<\alpha_0/\eta<4/7$ give
\begin{equation}\label{eq:decoder-b-D}
 \mathcal D(\{\infty\}\cup\mathcal V)
 \subset(-83/84,0),\qquad \frac5{12}+\frac47=\frac{83}{84}.
\end{equation}
This interval is disjoint from $\mathcal H$. Equations
\eqref{eq:decoder-a-E}--\eqref{eq:decoder-b-D} prove all of
\eqref{eq:decoder-wrong} for every nonempty directive.

\subsection{The empty directive and finite termination}

There is no last letter for the empty directive, so we check its
residual exclusions directly in the original coordinate. Here
$U=A$, $V=B$, $M_0=5$, and
$$
 K=\begin{pmatrix}3&4\\4&-3\end{pmatrix},\qquad
 E_0=B^{-1}KA=\begin{pmatrix}0&5\\5&-9\end{pmatrix},\qquad
 D_0=A^{-1}KB=\begin{pmatrix}9&5\\5&0\end{pmatrix}.
$$
The use of $K$ instead of $R$ does not change a projective map.
The terminal input is now two. We have $E_0(2)=5\notin \mathcal J$ and
$$
 E_0A(t)=\frac{5(t+1)}{t-4}<0\qquad(t\in\mathcal J).
$$
Thus an empty source remainder or a second $a$ is excluded.
For the other residual,
$$
 D_0(2)=23/10\in(2,50/21),\qquad
 D_0B(t)=\frac{55t+23}{25t+10}.
$$
The latter map is decreasing on $\mathcal J$ and satisfies
$$
 2<D_0B(t)\leqslant111/50<50/21.
$$
All these values lie between the terminal slope and the least
possible slope of a word starting with $B$. They lie outside the
encoded orbit. This proves exactly the empty-directive version
of \eqref{eq:decoder-wrong}, without a dominance assumption.

\begin{proof}[Proof of \cref{thm:word-decoder}]
Represent \eqref{eq:decoder-projective} in cusp coordinates.
If one word is empty, its slope is infinity; as $k(\infty)=\infty$
and no nonempty word has that slope, both words are empty.
Otherwise their first letters are opposite.

Suppose $x=ax_1$ and $y=by_1$. Removing the target's first
letter gives the residual relation
$y_1(\infty)=\mathcal E(x_1(\infty))$, where words here denote
their cusp maps. The first exclusion in
\eqref{eq:decoder-wrong} says that $x_1$ is neither empty nor
starts with $a$. Hence $x_1=bx_2$, and
\eqref{eq:decoder-residual-return} gives
$$
 y_1(\infty)=fk(x_2(\infty)).
$$
By \eqref{eq:decoder-positive-reflection} this is a finite negative
cusp slope. Therefore $y_1=ay_2$. Inverting $f$ restores
exactly the original relation
$y_2(\infty)=k(x_2(\infty))$.
Thus the first source block is $ab$ and the first target block
is $ba$.

If $x$ begins with $b$, the second exclusion in
\eqref{eq:decoder-wrong} and the identity $\mathcal Df=gk$
give the exchanged conclusion: the source begins with $ba$,
the target with $ab$, and removing those two letters restores
the same reflection relation. The empty-directive calculation
supplies the identical step in its original coordinate.

Each step removes two letters from each word. Induction on their
total length terminates, and the terminal case forces the two
words to end simultaneously. They are therefore precisely
the block-exchanged pair in \eqref{eq:decoder-projective}.
Conversely, repeated use of \eqref{eq:decoder-intertwiners} gives
$K\nu_d(x)=\nu_d(\widehat{x})K$ for every block word.
Applying both sides to $e$ proves \eqref{eq:decoder-exact} and
the converse implication.
\end{proof}

The exhaustion rests on the four strict residual exclusions and on the
termination of the two-letter descent, both established above. The
displayed identities and rational constants are also confirmed by exact
computation in the supplementary material; no enumeration of directives
or of finite words enters the proof.
